% Generated draft: requires mathematical and source review.
\documentclass[11pt]{article}
\usepackage[margin=25mm]{geometry}
\usepackage{fontspec}
\setmainfont{Latin Modern Roman}
\usepackage{amsmath,amssymb,mathtools,unicode-math}
\setmathfont{Latin Modern Math}
\usepackage{enumitem,needspace,xurl,hyperref}
\setlength{\parindent}{0pt}
\setlength{\parskip}{5pt}
\setlength{\emergencystretch}{3em}
\newcommand{\N}{\mathbb{N}}
\newcommand{\Z}{\mathbb{Z}}
\newcommand{\Q}{\mathbb{Q}}
\newcommand{\R}{\mathbb{R}}
\newcommand{\C}{\mathbb{C}}
\newcommand{\F}{\mathbb{F}}
\newcommand{\A}{\mathbb{A}}
\newcommand{\PP}{\mathbb{P}}
\newcommand{\E}{\mathbb{E}}
\newcommand{\eps}{\varepsilon}
\newcommand{\dd}{\,\mathrm{d}}
\begin{document}
\section*{UnsolvedMath ID: 700585}
\textbf{Source problem code:} catalogue entry
\subsection{Definitions and mathematical statement}
A group \(G\) is soluble if its derived series, defined by \(G^{(0)}=G\) and \(G^{(i+1)}=[G^{(i)},G^{(i)}]\), satisfies \(G^{(d)}=\{1\}\) for some integer \(d\geq 0\). It is finitely presented if it admits a presentation \(G=\langle X\mid R\rangle\) with both \(X\) and \(R\) finite. Interpret finite rank in the usual Pr\"ufer sense: there is an integer \(r\geq 0\) such that every finitely generated subgroup of \(G\) can be generated by at most \(r\) elements. The supplied record does not explicitly define its rank convention.

For a fixed finite presentation of \(G\), the conjugacy problem is decidable if there exists an algorithm which, on every pair of words \(u,v\) in the generators and their inverses, halts and decides whether
\[
\exists x\in G\qquad x^{-1}\overline{u}x=\overline{v},
\]
where \(\overline{u},\overline{v}\) are the elements represented by the words.

\textbf{Question.} Does every finitely presented soluble group of finite rank have a decidable conjugacy problem? Here ``soluble conjugacy problem'' is read as algorithmic decidability. The formulation asks for an algorithm for each such group; the supplied record does not specify a uniform algorithm taking arbitrary presentations as additional input.

\textbf{Supplied qualification.} The source adds that there is an algorithm to decide conjugacy to a given element of the group. Read as a fixed-target assertion, this says that for each fixed \(g\in G\), there exists an algorithm deciding, for every input word \(u\), whether \(\overline{u}\) is conjugate to \(g\). The record does not specify effective uniformity in \(g\); its note must therefore be distinguished from an algorithm receiving both elements as inputs.
\subsection{Short English statement}
Must every finitely presented soluble group of finite Pr\"ufer rank admit an algorithm deciding whether any two input words represent conjugate elements? The supplied note asserts decidability for a fixed target element, without specifying uniformity in that target. The rank convention and algorithmic interpretation require human review.
\subsection{Sources}
\begin{itemize}
\item Local UnsolvedMath registry, permanent ID \(700585\), exact source problem code \texttt{catalogue entry}, title ``Kourovka Notebook Problem 9.52.'' The supplied registry labels the problem open; that label is supplied curation, not an independently verified mathematical or literature assessment. The record was created and updated on 2026-10-16.
\item Attributed in the supplied record to D. J. S. Robinson, proposed in 1984, as Problem 9.52 of \emph{The Kourovka Notebook (Unsolved Problems in Group Theory)}, 9th issue (1984). The record also identifies the 21st issue edition (2026), arXiv:1401.0300v48. Its collection description names E. I. Khukhro and V. D. Mazurov as editors. These bibliographic details are supplied attribution and have not been independently checked.
\end{itemize}
\end{document}
